% GENERATED FILE. Edit canonical lessons, then run npm run generate:books.
% canonical-source-hash: fnv1a64:238e4695cb06ac42
% canonical-lessons: ML-C343-rodu, ML-W343-da, ML-C343-evideyanu, ML-R343-asking-the-way

\chapter{Asking the Way}
\label{ch:ml-asking-the-way}
\hlchaptermodality{343}

\begin{chapteropening}
\textbf{By the end of this chapter:} \emph{I can ask where a place is and how to get there, and follow and give simple directions: straight on, left, right, at the junction.}

Complete ML-R343-asking-the-way: ask for the station, follow and give a route with nēre, iṭattōṭṭŭ and valattōṭṭŭ, and read the new letter \ml{ഡ} in \ml{റോഡ്}.
\end{chapteropening}

\section[\texorpdfstring{\ml{റോഡ്} (ṟōḍŭ)}{റോഡ് (ṟōḍŭ)}]{\ml{റോഡ്} (ṟōḍŭ) — a road}

\label{lesson:ML-C343-rodu}



\begin{quote}
Before the new one: say the Malayalam for \textquotedblleft{}a way, a path\textquotedblright{}, then \textquotedblleft{}a street\textquotedblright{}, then \textquotedblleft{}a station\textquotedblright{}. (\emph{vaḻi}, \emph{teruvŭ}, \emph{sṟṟēṣan}.)
\end{quote}

\subsection*{You'll want to know: \ml{റോഡ്}}
\textbf{\ml{റോഡ്}} — \emph{ṟōḍŭ} — \textquotedblleft{}a road\textquotedblright{}.

\begin{quote}
\textbf{\ml{ഈ} \ml{റോഡ്} \ml{സ്റ്റേഷനിലേക്ക്} \ml{പോകുന്നു}.}
\end{quote}

\begin{quote}
\emph{ī ṟōḍŭ sṟṟēṣanilēkkŭ pōkunnu.} — \textquotedblleft{}This road goes to the station.\textquotedblright{}
\end{quote}

\begin{quote}
\textbf{\ml{വലിയ} \ml{റോഡിൽ} \ml{ഒരുപാട്} \ml{ബസുകൾ} \ml{ഉണ്ട്}.}
\end{quote}

\begin{quote}
\emph{valiya ṟōḍil orupāṭŭ basukaḷ uṇṭŭ.} — \textquotedblleft{}There are a lot of buses on the main road.\textquotedblright{}
\end{quote}

\textbf{\ml{വഴി}} is the way or route you take; \textbf{\ml{റോഡ്}} is the road itself; \textbf{\ml{തെരുവ്}} is a street lined with houses and shops. \textbf{-\ml{ലേക്ക്}} on \textbf{\ml{സ്റ്റേഷൻ}} says \textquotedblleft{}towards, to\textquotedblright{}.

Its last consonant is new. The next lesson takes it out of this word and puts it in your hand.

\begin{cousinweb}
\textbf{\ml{റോഡ്}} (\emph{ṟōḍŭ}) is borrowed from English \emph{road}. Malayalam hears the English \emph{r} as its own \textbf{\ml{റ}} and the English \emph{d} as its curled-back \textbf{\ml{ഡ}}, and closes the word with the half-u, \emph{ŭ}, as it does with \textbf{\ml{ടിക്കറ്റ്}} and \textbf{\ml{സീറ്റ്}}.
\end{cousinweb}

\subsection*{Your turn}
Say these aloud:

\begin{itemize}
\raggedright
  \item \emph{ṟōḍŭ}, with the tongue curled back for the \emph{ḍ}
  \item \textquotedblleft{}This road goes to the station.\textquotedblright{}
  \item \textquotedblleft{}There are a lot of buses on the main road.\textquotedblright{}
  \item a way, a road, a street — \emph{vaḻi}, \emph{ṟōḍŭ}, \emph{teruvŭ}
\end{itemize}

\begin{tcolorbox}[breakable,colback=teal!4,colframe=teal!35!black,title={Before you move on}]
What does \textbf{\ml{റോഡ്}} mean? (\textbf{a road.}) How is it different from \textbf{\ml{വഴി}}? (\textbf{\ml{വഴി} is the way you take; \ml{റോഡ്} is the road itself.})
\end{tcolorbox}
\section[\texorpdfstring{\ml{ഡ} (ḍa)}{ഡ (ḍa)}]{\ml{ഡ} (ḍa) — the letter inside \ml{റോഡ്}}

\label{lesson:ML-W343-da}



\begin{quote}
Before the new one: write \textbf{\ml{ട}} and \textbf{\ml{ദ}} once each, and name them. (\emph{ṭa}, with the tongue curled back; \emph{da}, with the voice on.)

Now say \textbf{\ml{റോഡ്}} (\emph{ṟōḍŭ}), \textquotedblleft{}a road\textquotedblright{}. This lesson takes its last consonant out and puts it in your hand.
\end{quote}

\subsection*{Script you'll notice: \ml{ഡ}}
\textbf{\ml{ഡ}} — \emph{ḍa}: the tongue curled back as for \textbf{\ml{ട}}, and the voice on as for \textbf{\ml{ദ}}. It is the voiced partner of \textbf{\ml{ട}}.

In \textbf{\ml{റോഡ്}} it carries the half-u sign \textbf{\ml{്}}, so the word ends \emph{-ḍŭ}. Native Malayalam words rarely use it; it is the letter Malayalam reaches for when an English word has a \emph{d}.

\subsection*{Writing: \ml{ഡ}}
Put your pen on \textbf{\ml{ഡ}} and follow its line. Copy the shape you can see, slowly, and larger than it is printed.

\begin{quote}
This book does not yet tell you \textbf{where to start this letter or which way to travel}. It is not written down here until it can be written down with a source. Copying what is in front of you needs no such source, and it is how the shape gets into your hand in the meantime.
\end{quote}

\subsection*{Your turn}
\begin{itemize}
\raggedright
  \item \emph{Find:} \ml{ഡ} inside \textbf{\ml{റോഡ്}}
  \item \emph{Trace it:} \ml{ഡ} three times, saying \emph{ḍa} as you finish each one
  \item \emph{Write it:} \ml{ട}, then \ml{ഡ}, from memory
  \item \emph{Write it:} \textbf{\ml{റോഡ്}}
\end{itemize}

\begin{tcolorbox}[breakable,colback=teal!4,colframe=teal!35!black,title={Before you move on}]
Which letter is this — \ml{ഡ}? (\textbf{ḍa}: curled back, voiced.) Which word did it come from? (\textbf{\ml{റോഡ്}}, \emph{ṟōḍŭ}, \textquotedblleft{}a road\textquotedblright{}.) Which letter is its unvoiced partner? (\textbf{\ml{ട}}.)
\end{tcolorbox}
\section[\texorpdfstring{\ml{സ്റ്റേഷൻ} \ml{എവിടെയാണ്}? (sṟṟēṣan eviṭeyāṇŭ?)}{സ്റ്റേഷൻ എവിടെയാണ്? (sṟṟēṣan eviṭeyāṇŭ?)}]{\ml{സ്റ്റേഷൻ} \ml{എവിടെയാണ്}? (sṟṟēṣan eviṭeyāṇŭ?) — where is the station?}

\label{lesson:ML-C343-evideyanu}



\begin{quote}
Say \textquotedblleft{}where?\textquotedblright{}, then \textquotedblleft{}on the left\textquotedblright{}, \textquotedblleft{}on the right\textquotedblright{}, and \textquotedblleft{}straight ahead\textquotedblright{}. (\emph{eviṭe}, \emph{iṭattŭ}, \emph{valattŭ}, \emph{nēre}.)
\end{quote}

\begin{grammarlens}[title={Grammar Lens: asking}]
Open with \textbf{\ml{ക്ഷമിക്കണം}}, \textquotedblleft{}excuse me\textquotedblright{}, then ask where, or ask how:

\begin{quote}
\textbf{\ml{ക്ഷമിക്കണം}, \ml{സ്റ്റേഷൻ} \ml{എവിടെയാണ്}?}
\end{quote}

\begin{quote}
\emph{kṣamikkaṇaṁ, sṟṟēṣan eviṭeyāṇŭ?} — \textquotedblleft{}Excuse me, where is the station?\textquotedblright{}
\end{quote}

\begin{quote}
\textbf{\ml{സ്റ്റേഷനിലേക്ക്} \ml{എങ്ങനെ} \ml{പോകണം}?}
\end{quote}

\begin{quote}
\emph{sṟṟēṣanilēkkŭ eṅṅane pōkaṇaṁ?} — \textquotedblleft{}How do I get to the station?\textquotedblright{}
\end{quote}

The second is the \textbf{-\ml{അണം}} of chapter 342: \textquotedblleft{}how must one go\textquotedblright{}. \textbf{\ml{എത്ര} \ml{ദൂരമുണ്ട്}}, \emph{etra dūramuṇṭŭ}, asks how far.
\end{grammarlens}

\begin{grammarlens}[title={Grammar Lens: following the answer}]
\textbf{\ml{ഇടത്ത്}} is \textquotedblleft{}on the left\textquotedblright{}; add \textbf{-\ml{ഓട്ട്}}, as in \textbf{\ml{അങ്ങോട്ട്}}, and it is \textquotedblleft{}to the left\textquotedblright{}: \textbf{\ml{ഇടത്തോട്ട്}}, \emph{iṭattōṭṭŭ}. So \textbf{\ml{വലത്തോട്ട്}}, \textquotedblleft{}to the right\textquotedblright{}.

\begin{quote}
\textbf{\ml{ഈ} \ml{റോഡിലൂടെ} \ml{നേരെ} \ml{പോകൂ}. \ml{കവലയിൽ} \ml{എത്തുമ്പോൾ} \ml{ഇടത്തോട്ട്} \ml{തിരിയൂ}. \ml{സ്റ്റേഷൻ} \ml{വലത്ത്} \ml{ആണ്}.}
\end{quote}

\begin{quote}
\emph{ī ṟōḍilūṭe nēre pōkū. kavalayil ettumpōḷ iṭattōṭṭŭ tiriyū. sṟṟēṣan valattŭ āṇŭ.} — \textquotedblleft{}Go straight along this road. When you reach the junction, turn left. The station is on the right.\textquotedblright{}
\end{quote}

\textbf{-\ml{ലൂടെ}} is \textquotedblleft{}along, through\textquotedblright{}, and \textbf{\ml{രണ്ടാമത്തെ} \ml{റോഡ്}} is \textquotedblleft{}the second road\textquotedblright{}.
\end{grammarlens}

\subsection*{Your turn}
Say these aloud:

\begin{itemize}
\raggedright
  \item \textquotedblleft{}Excuse me, where is the station?\textquotedblright{}
  \item \textquotedblleft{}How do I get to the station?\textquotedblright{}, then ask how far it is
  \item \textquotedblleft{}to the left\textquotedblright{}, \textquotedblleft{}to the right\textquotedblright{} — \emph{iṭattōṭṭŭ}, \emph{valattōṭṭŭ}
  \item the route: straight along this road, left at the junction, the station on the right
\end{itemize}

\begin{tcolorbox}[breakable,colback=teal!4,colframe=teal!35!black,title={Before you move on}]
\textquotedblleft{}How do I get to the station?\textquotedblright{} (\emph{\textbf{sṟṟēṣanilēkkŭ eṅṅane pōkaṇaṁ?}}) \textquotedblleft{}Turn left\textquotedblright{}? (\emph{\textbf{iṭattōṭṭŭ tiriyū.}}) \textquotedblleft{}Along this road\textquotedblright{}? (\emph{\textbf{ī ṟōḍilūṭe.}})
\end{tcolorbox}
\section[(review)]{Review — asking the way}

\label{lesson:ML-R343-asking-the-way}



\begin{quote}
Say \textquotedblleft{}a road\textquotedblright{}. (\emph{ṟōḍŭ}.) Write \textbf{\ml{ഡ}} once from memory.
\end{quote}

\subsection*{Your turn}
\begin{itemize}
\raggedright
  \item \emph{Say it:} \textquotedblleft{}Excuse me, where is the station?\textquotedblright{}
  \item \emph{Say it:} \textquotedblleft{}How do I get to the station? How far is it from here?\textquotedblright{}
  \item \emph{Say it:} \textquotedblleft{}Go straight along this road. At the junction turn right.\textquotedblright{}
  \item \emph{Say it:} \textquotedblleft{}Because the road is busy, take the second road.\textquotedblright{} — \emph{ṟōḍil tirakkuḷḷatukoṇṭŭ raṇṭāmatte ṟōḍil pōkū.}
  \item \emph{Say it:} the way from your door to the nearest shop, in three steps
  \item \emph{Say it:} \textquotedblleft{}Tomorrow I will go to the station, because I want a ticket.\textquotedblright{} — \emph{nāḷe ñān sṟṟēṣanil pōkuṁ, kāraṇaṁ enikkŭ oru ṭikkaṟṟŭ vēṇaṁ.}
  \item \emph{Write it:} \textbf{\ml{റോഡ്}}
\end{itemize}

\begin{tcolorbox}[breakable,colback=teal!4,colframe=teal!35!black,title={Before you move on}]
\textquotedblleft{}Where is …?\textquotedblright{} (\textbf{… \ml{എവിടെയാണ്}?}) \textquotedblleft{}How do I get to …?\textquotedblright{} (\textbf{… -\ml{ലേക്ക്} \ml{എങ്ങനെ} \ml{പോകണം}?}) \textquotedblleft{}To the right\textquotedblright{}? (\textbf{\ml{വലത്തോട്ട്}}.) \textquotedblleft{}A road\textquotedblright{}? (\textbf{\ml{റോഡ്}}, which ends in \textbf{\ml{ഡ്}}.)
\end{tcolorbox}
